\section{Các công trình liên quan}
\subsection{Mô hình CNN nhẹ cho phân đoạn ngữ nghĩa}
Semantic segmentation đã đạt nhiều tiến bộ nhờ các mạng CNN sâu, bắt đầu từ hướng dự đoán dày đặc ở mức điểm ảnh của FCN \cite{long2015fcn}. Tuy nhiên, khi triển khai trên thiết bị có tài nguyên hạn chế, các mô hình lớn thường gặp rào cản về độ trễ suy luận, bộ nhớ và chi phí tính toán. Vì vậy, một nhánh nghiên cứu quan trọng tập trung vào \emph{lightweight CNN}, hướng tới duy trì độ chính xác cạnh tranh trong khi giảm số tham số và FLOPs.

Các mô hình lightweight tiêu biểu theo đuổi nhiều chiến lược khác nhau. ENet mở đầu xu hướng thời gian thực với thiết kế encoder--decoder tối giản và bottleneck hiệu quả \cite{paszke2016enet}. BiSeNet \cite{bisenet} tiếp tục cải thiện bài toán cân bằng tốc độ--độ chính xác bằng kiến trúc hai nhánh (bilateral): nhánh không gian bảo toàn chi tiết biên và nhánh ngữ cảnh mở rộng trường thụ cảm; nhờ đó mô hình có thể đạt tốc độ cao (ví dụ 105 FPS trên Cityscapes) mà vẫn giữ chất lượng phân đoạn tốt. Bên cạnh đó, nhiều kiến trúc nhẹ khác như LEDNet, CFPNet hay LETNet khai thác các kỹ thuật như depthwise separable convolution, channel shuffling và thiết kế feature pyramid gọn hơn. Dù hiệu quả ở các cảnh tổng quát, nhóm phương pháp này vẫn có thể hạn chế khi cần mô hình hóa phụ thuộc tầm xa trong các ngữ cảnh phức tạp và chi tiết mịn.

\subsection{CCNet và cơ chế chú ý Criss-Cross}
Trên nền tảng các backbone CNN hiệu quả, CCNet \cite{ccnet} là một bước tiến đáng chú ý trong mô hình hóa ngữ cảnh cho phân đoạn ngữ nghĩa. Thay vì dùng attention toàn cục kiểu non-local với độ phức tạp cao, CCNet đề xuất \emph{criss-cross attention} để mỗi pixel tổng hợp thông tin theo hàng và cột tương ứng. Cơ chế này giảm đáng kể chi phí so với attention toàn cặp nhưng vẫn duy trì khả năng nắm bắt ngữ cảnh xa.

Để đạt phụ thuộc toàn ảnh, CCNet áp dụng \emph{Recurrent Criss-Cross Attention} (RCCA) theo hướng lặp (thường dùng $R=2$), giúp thông tin lan truyền qua nhiều bước và bao phủ các vùng xa trong ảnh. Mô hình cũng sử dụng \emph{category consistent loss} để tăng tính phân biệt giữa các lớp ngữ nghĩa. Thực nghiệm trong bài gốc cho thấy CCNet đạt hiệu năng mạnh trên các benchmark lớn (như Cityscapes, ADE20K, PASCAL VOC/LIP tùy thiết lập) với chi phí thấp hơn đáng kể so với các phương pháp attention toàn cục trước đó \cite{ccnet}.

\subsection{CCNet trong phân đoạn ảnh thực phẩm}
Trong miền \emph{food image segmentation}, bộ dữ liệu FoodSeg103 được giới thiệu như một benchmark quy mô lớn và nhiều thách thức, với gán nhãn pixel-wise ở mức thành phần thực phẩm \cite{benchmarkfood}. Các khó khăn đặc thù của miền này gồm biến thiên nội lớp lớn, che khuất và tương đồng mạnh về màu sắc/kết cấu giữa các món, khiến nhu cầu mô hình hóa ngữ cảnh tầm xa trở nên quan trọng.

Trên thiết lập benchmark của FoodSeg103, CCNet (ResNet-50) được sử dụng như một baseline convolutional mạnh và cho kết quả cạnh tranh; trong báo cáo gốc, mIoU được nêu khoảng $35.5\%$ cho CCNet chuẩn và tăng lên khoảng $36.8\%$ khi kết hợp mô-đun tinh chỉnh ReLeM \cite{benchmarkfood}. Kết quả này cho thấy cơ chế attention có mục tiêu như RCCA có thể cải thiện khả năng phân đoạn thành phần thực phẩm mà vẫn giữ tính khả thi triển khai. Tổng thể, tiến trình từ CNN nhẹ đến attention hiệu quả như CCNet cho thấy hướng đi rõ ràng cho các hệ thống semantic segmentation thực tiễn: cân bằng chặt chẽ giữa hiệu quả tính toán và năng lực thu ngữ cảnh giàu thông tin.